% !TeX root = ../MAgent_Theory_2.0_companion.tex
\chapter{What a Message Lets the Receiver Infer}
\label{comp:communication}

Sending a few observables can change an agent's description of many variables. The extra structure comes from the receiver's conditional model. This is familiar in physics: measuring one coordinate of a correlated system can change predictions for another coordinate without measuring it directly. The distinction between transmitted information and inferred structure is explicit here.

\section{Conditional prediction within an agent}
\label{comp:conditional-prediction}

An agent can use its model to infer a result from a different, correlated cue. Let $A\in\{1,\ldots,M\}$ be a finite model presentation with prior weights $\pi_a>0$ summing to one. On standard Borel cue and result spaces, presentation $a$ supplies a cue law $\mu_a(du)$ and a normalized result kernel $K_a(u,dv)$. A normalized record kernel $\mathcal L(u,do)$ measures the cue. The joint law is
\begin{equation}
P(a,du,dv,do)=\pi_a\mu_a(du)K_a(u,dv)\mathcal L(u,do).
\label{comp:eq-cue-result-joint}
\end{equation}
Successive integration of the kernels and summation over $a$ give total mass one. In the agent description, $a$ is the model presentation, $(u,v)$ is the represented content, and $\mu_a(du)K_a(u,dv)$ is its evaluator. The kernels and prior weights constitute the stored model; the posterior after a record describes the agent's current beliefs.

Suppose $\mathcal L(u,do)=l(o\mid u)\lambda(do)$ and the observed $o$ has finite positive evidence $Z_a(o)=\int l(o\mid u)\mu_a(du)$ in every component. Bayes' rule gives
\begin{equation}
\begin{aligned}
w_a&=\frac{\pi_aZ_a(o)}{\sum_b\pi_bZ_b(o)},\\
\Pi_a^o(du,dv)&=\frac{l(o\mid u)\mu_a(du)K_a(u,dv)}{Z_a(o)}.
\end{aligned}
\label{comp:eq-cue-result-posterior}
\end{equation}
The exact recognition law is $w_a\Pi_a^o$, so the model belief $s(a)=w_a$ is generally correlated with the contents. A product restriction would replace the component-dependent content laws by one common law. Conditioning revises these beliefs without changing the stored kernels. This construction uses general probability laws; Gaussian conditioning makes a particular realization explicit.

For a scalar example, choose one presentation and independent variables
\begin{equation}
X\sim\mathcal N(0,1),\qquad
\epsilon\sim\mathcal N(0,1/4),\qquad
\eta\sim\mathcal N(0,1),
\end{equation}
with result $Y=-2X+\epsilon$ and record $O=X+\eta$. Observe $O=\xi$. The cue posterior follows by completing the square:
\begin{equation}
p(x\mid\xi)\propto
\exp\left[-\frac{x^2+(\xi-x)^2}{2}\right]
\propto\exp[-(x-\xi/2)^2],
\end{equation}
hence $X\mid O=\xi\sim\mathcal N(\xi/2,1/2)$. Since the result noise is independent of the cue and record,
\begin{equation}
Y\mid O=\xi\sim\mathcal N(-\xi,9/4).
\label{comp:eq-negative-result}
\end{equation}
The variance is $4(1/2)+1/4=9/4$. Increasing the observed cue increases the inferred cue but decreases the predicted result. The negative response comes from the model's conditional relation between two variables. Estimating the cue itself in this example would instead have positive slope $1/2$. Conditional prediction therefore gives a statistical account of a transformation that reverses a direction. Conditioning a joint Gaussian mixture to predict one group of variables from another is Gaussian mixture regression \citep{Calinon2016}.

With several presentations, uncertainty also concerns which relation applies. For vector-valued results with finite second moments, let $m_a=\E[Y\mid o,A=a]$ and $S_a=\operatorname{Cov}(Y\mid o,A=a)$. The result marginal $p_Y^o$ has moments
\begin{equation}
\begin{aligned}
\bar m&=\sum_a w_am_a,\\
S&=\sum_a w_a\left[S_a+(m_a-\bar m)(m_a-\bar m)^\top\right].
\end{aligned}
\label{comp:eq-result-ambiguity}
\end{equation}
The first covariance contribution is uncertainty within a presentation. The second is uncertainty from the disagreement of its predictions with those of other presentations. These terms follow by writing $Y-\bar m=(Y-m_a)+(m_a-\bar m)$ and taking the conditional expectation; the cross terms vanish.

If the result is a finite mixture of nonsingular Gaussians, one possible carrier is $q_G=\mathcal N(\bar m,S)$. It minimizes $\KL(p_Y^o\Vert q_G)$ over Gaussian laws. Indeed, the expected negative Gaussian log density depends on $p_Y^o$ only through its first two moments, and its minimum occurs at their matching mean and covariance. This is the moment projection used in assumed-density filtering \citep{Minka2001EP}. The variational free energy instead contains the divergence with recognition in the first slot. Minimizing $\KL(q\Vert p_Y^o)$ over Gaussians is a different problem and need not return the moment carrier.

Retaining $w$ alongside $q_G$ still does not preserve their original dependence. If the retained description is the product $w\otimes q_G$, the exact label/result posterior $\Pi_{A,Y}^o$ obeys
\begin{equation}
\KL(\Pi_{A,Y}^o\Vert w\otimes q_G)
=I_{\Pi^o}(A;Y)+\KL(p_Y^o\Vert q_G).
\label{comp:eq-memory-retention-loss}
\end{equation}
Here $I_{\Pi^o}(A;Y)=\KL(\Pi_{A,Y}^o\Vert w\otimes p_Y^o)$ is mutual information. Insert $w\otimes p_Y^o$ into the logarithmic density ratio to obtain the identity. It separates the loss of model/content dependence from the loss of non-Gaussian result information. Neither loss is determined by the accuracy of the mean alone.

A later step may carry the component belief forward, for example as $wT$ under a declared stochastic transition matrix $T$. An interpretation as Bayesian filtering also needs the conditional law of the next record given the retained history and presentation. A deterministic function of a record already available to the agent supplies no further evidence conditional on that record; treating it as an independent observation would change the model. Communication can supply information new to the receiver, but its record kernel must be specified. The local conditioning construction therefore provides an agent computation, while repeated assimilation, model learning, and population dynamics require their own declared mechanisms.

\section{Selective sending and conditional reconstruction}

A sender uses a measurable map $T_j:E_j\to\mathsf H$ into a message space. If its content has law $q_j$, possible messages have law
\begin{equation}
\mu_j=(T_j)_\#q_j.
\label{comp:eq-message-law}
\end{equation}
The pushforward symbol means $\mu_j(A)=q_j(T_j^{-1}A)$. Sampling a message and transmitting the whole law $\mu_j$ are different communication protocols. Messages run from sender to receiver. The reciprocal potential of the dynamics chapter also moves each sender toward its listeners, which a communicating population can only realize if the listeners send something back. A linear receiving map $O_i:\mathsf H\to E_i$ embeds a message, but its composition with $T_j$ cannot have higher rank than $T_j$.

\begin{figure}[htbp]
\centering
\begin{tikzpicture}[
  >={Stealth[length=2.2mm]},
  box/.style={draw=primaryblue,rounded corners=2pt,fill=primaryblue!5,minimum width=2.55cm,minimum height=1.05cm,align=center,font=\small},
  every node/.style={font=\small}]
\node[box] (source) {Sender law\\$q(dk)$};
\node[box,right=1.6cm of source] (message) {Message law\\$\nu=T_\#q$};
\node[box,right=1.6cm of message] (receiver) {Reconstruction\\$\widehat q(dk)$};
\node[box,below=1.1cm of receiver,fill=secondaryblue!10] (prior) {Receiver conditional\\$p(dk\mid h)$};
\draw[->,thick,primaryblue] (source) -- node[above=4pt] {$T$} (message);
\draw[->,thick,primaryblue] (message) -- node[above=4pt,align=center] {mix over $h$} (receiver);
\draw[->,thick,secondaryblue] (prior) -- (receiver);
\node[below=0.65cm of source,align=center,text width=3cm] {Different sender laws\\can give the same $\nu$.};
\end{tikzpicture}

\caption{Selective communication and receiver structure. The sender fixes the transmitted marginal. The receiver supplies a conditional law for distinctions that the message does not specify. This distributional reconstruction presupposes that the message law is available; one received sample supplies a value instead.}
\label{fig:comp-message}
\end{figure}

To see exactly what is lost, work in one common presentation and write $T:E\to\mathsf H$. Let $q$ be the sender's law and $p$ the receiver's reference law, with finite $\KL(q\Vert p)$. The receiver can combine the transmitted marginal $\mu=T_\#q$ with its own conditional law:
\begin{equation}
\widehat q(dk)=\int p(dk\mid h)\mu(dh).
\label{comp:eq-reconstruction}
\end{equation}
Disintegration on standard Borel spaces provides these conditional laws \citep{Kallenberg2021}. The reconstruction has exactly the right message marginal. Its remaining discrepancy is
\begin{equation}
\KL(q\Vert\widehat q)
=\int\KL\bigl(q(\cdot\mid h)\Vert p(\cdot\mid h)\bigr)\mu(dh).
\label{comp:eq-reconstruction-discrepancy}
\end{equation}
Both laws assign the same probability to $h$, so that factor cancels in their density ratio. The remaining logarithm compares the two conditional laws, which gives the equality by integration. This reconstruction is Jeffrey's rule of conditioning \citep{DiaconisZabell1982}, and it is the law closest to the receiver's prior among all laws with the transmitted marginal. The divergence then splits like a right triangle, $\KL(q\Vert p)=\KL(q\Vert\widehat q)+\KL(\widehat q\Vert p)$, with the second leg equal to the divergence of the message marginals \citep{Csiszar1975}.

For a calculation on four points, let $k=(H,V)$ be two bits. The sender gives probability $1/2$ to each of $00$ and $11$, and transmits only $H$. A receiver whose reference makes both bits independent and fair reconstructs the uniform law on all four points. The message marginal is correct, while the relative entropy is $\log2$. A receiver whose conditional model has $V=H$ reconstructs the sender's full law. The missing correlation was supplied by the receiver.

One reference conditional can recover a whole family only under a stronger condition. Suppose $q_\alpha=L_\alpha p$ and use the same $p(dk\mid h)$ for every $\alpha$. Its reconstruction has density
\begin{equation}
\frac{d\widehat q_\alpha}{dp}(k)
=\E_p[L_\alpha\mid T](k).
\end{equation}
To obtain this identity, weight the reference message marginal by its conditional mean likelihood and then apply the reference conditional. The result equals $q_\alpha$ precisely when $L_\alpha$ is already measurable from the message, which is the classical statement that the message is a sufficient statistic for the family \citep{HalmosSavage1949}. Reconstructing one posterior with its own conditional does not establish this family-wide property. If a local reconstruction must also ignore retained exterior information, the reference needs conditional independence from that exterior given the message.

For a single received value $h$, Bayesian assimilation instead takes the form
\begin{equation}
q_i^+(dk)=\frac{\ell_i(h\mid k)q_i^-(dk)}{\int\ell_i(h\mid k')q_i^-(dk')},
\label{comp:eq-bayes-receiver}
\end{equation}
with a positive finite denominator. In the perfectly correlated two-bit example, a noiseless observation $H=1$ also fixes $V=1$. This is inference through an existing relation, not an inversion of an arbitrary lossy map. Under the reciprocal potential such an update is a kick: the flow afterwards relaxes toward states that do not depend on the record, unless the record's likelihood is kept in the objective.

\section{Agreement needs a common prediction rule}

Suppose a model law $s$ and a normalized record evaluator $K(m,do)$ predict the record law $sK$. If two agents compare model laws in one frame, they must also compare their evaluators in that frame. Equal model laws can predict opposite bits when their evaluators differ.

For one common evaluator and finite $\KL(s\Vert t)$, an exact information balance is
\begin{equation}
\KL(s\Vert t)
=\KL(sK\Vert tK)
+\int\KL(s^o\Vert t^o)(sK)(do),
\label{comp:eq-prediction-balance}
\end{equation}
where $s^o,t^o$ are the conditional model laws after record $o$. The joint laws $s(dm)K(m,do)$ and $t(dm)K(m,do)$ have divergence $\KL(s\Vert t)$ because their common kernel cancels. Factoring the same joint laws by their record marginal gives the displayed split. The last term is the distinction left unresolved by the record. A constant evaluator discards every distinction; a sufficient record leaves none.

For a model link $\Theta^s_{ij}$, compatibility means $K_i(\Theta^s_{ij}m,A)=K_j(m,A)$ after identifying record alphabets. Only then does transporting the model law also transport its predictive meaning. If the link changes with the laws, this condition must continue to hold at the updated state.

\section{Compression should name what matters}

Rate--distortion theory makes the desired fidelity part of the communication problem \citep{Shannon1959}. A low-dimensional real vector alone is not a finite-bit channel. Noise, precision, or an encoding constraint must specify the communication cost. Wyner--Ziv coding additionally gives the receiver a realized variable correlated with the source \citep{WynerZiv1976}; possession of a prior by itself is not that realized side information.

For a predictive criterion, name the source $X$, target response $Y$, and accessible receiver information $Z$. An encoder generates $U$ from $X$ alone, so $U-X-(Y,Z)$ is a conditional-independence chain. The optimal decoder is $P(Y\mid U,Z)$. Its expected relative-entropy distortion is
\begin{equation}
\mathcal D=I(Y;X\mid U,Z)
=I(Y;X\mid Z)-I(Y;U\mid Z).
\label{comp:eq-predictive-distortion}
\end{equation}
The equality follows by expanding $I(Y;X,U\mid Z)$ in both orders. It is the information about the chosen response that compression discards. An information-bottleneck criterion trades this quantity against a cost such as $I(X;U\mid Z)$ \citep{Tishby1999}; an operational bit rate requires additional coding assumptions.

If $X=(A,B)$ consists of independent fair bits and $U=A$, the distortion is zero for target $Y=A$ and $\log2$ for target $Y=B$. For a prospective parent, the target can be a future boundary response. Retaining one such response does not ensure that every later interaction can be inherited. That stronger requirement concerns dynamics.
